\documentclass[11pt]{article}

\usepackage[a4paper, margin=27mm]{geometry}
\usepackage{fontspec}
\setmainfont{Times New Roman}[   % Times-like, full Vietnamese coverage
  Path = /System/Library/Fonts/Supplemental/,
  Extension = .ttf,
  UprightFont = *,
  BoldFont = * Bold,
  ItalicFont = * Italic,
  BoldItalicFont = * Bold Italic,
]
\usepackage{microtype}
\usepackage{booktabs}
\usepackage{amsmath}
\usepackage{enumitem}
\usepackage{listings}
\usepackage{xcolor}
\usepackage[hidelinks]{hyperref}
\usepackage{parskip}

\setlength{\parskip}{4pt plus 1pt}
\setlength{\parindent}{0pt}

\lstset{
  basicstyle=\footnotesize\ttfamily,
  frame=single,
  framerule=0.3pt,
  rulecolor=\color{black!40},
  backgroundcolor=\color{black!4},
  xleftmargin=8pt, xrightmargin=8pt,
  aboveskip=8pt, belowskip=8pt,
  breaklines=true,
  columns=fullflexible,
  keepspaces=true,
}

% Listings wrapped in a minipage cannot break across pages
\lstnewenvironment{code}
  {\medskip\noindent\minipage{\linewidth}}
  {\endminipage\medskip}

\begin{document}

\begin{center}
{\Large\bfseries Freeport: Giao thức cho một Chợ Ngang hàng}\\[10pt]
{\normalsize Phil Trinh}\\[2pt]
{\small\texttt{ptrinh@me.com}}\\[8pt]
{\small Bản thảo v0.2 --- Tháng 6, 2026}
\end{center}

\vspace{14pt}

\begin{center}
\begin{minipage}{0.88\textwidth}
\small
\textbf{Tóm tắt.} Một chợ thuần ngang hàng sẽ cho phép người mua và người
bán hàng hóa, dịch vụ địa phương giao dịch trực tiếp với nhau mà không phải
đi qua một nền tảng trung gian. Chữ ký số và các relay mở giải quyết được một
phần bài toán, nhưng phần lớn lợi ích sẽ mất đi nếu vẫn cần một bên thứ ba
đáng tin để thiết lập uy tín. Chúng tôi đề xuất một lời giải cho bài toán
niềm tin bằng một mạng lưới biên nhận giao dịch được đồng ký. Người tham gia
công bố các \emph{ý định} (intent) đã ký lên các relay mở; các bên đối tác
thương lượng điều khoản qua tin nhắn mã hóa bằng một máy trạng thái có giới
hạn; một giao dịch chỉ được \emph{chứng minh} khi cả hai bên độc lập ký một
biên nhận tham chiếu đến nó. Uy tín không phải là một điểm số toàn cục mà là
một đại lượng chủ quan: mỗi người xem tự tính bằng cách lần theo đồ thị biên
nhận từ chính khóa của mình, nên các danh tính giả mạo --- dù đông đến đâu ---
cũng không có trọng lượng với những ai họ chưa thực sự phục vụ. Giao thức
không cần nền tảng, không phí, không tài khoản và không cần xin phép. Các chợ
được định nghĩa bằng quy ước: bất kỳ nhóm người nào thống nhất một thẻ chủ đề
(topic tag) và một lược đồ dữ liệu (payload schema) đều tạo thành một chợ.
\end{minipage}
\end{center}

\section{Giới thiệu}

Thương mại giữa các cá nhân trên internet đã gần như phụ thuộc hoàn toàn vào
các nền tảng đóng vai trò bên thứ ba đáng tin: các công ty gọi xe, các chợ
việc làm tự do, các trang rao vặt. Tuy mô hình này hoạt động đủ tốt cho phần
lớn giao dịch, nó mang đầy đủ những điểm yếu cố hữu của mô hình dựa trên niềm
tin. Các nền tảng thu hoa hồng 15--30\% giá trị giao dịch. Họ đơn phương đặt
ra điều khoản, khóa tài khoản mà không cho khiếu nại, và giữ uy tín làm con
tin --- nhiều năm đánh giá của một người bán không thể rời khỏi nền tảng đang
lưu trữ chúng. Hiệu ứng mạng củng cố vị thế của kẻ đi trước; cả hai phía đều
phải trả tô độc quyền. Nền tảng phân xử các tranh chấp mà chính nó có động cơ
giải quyết theo hướng có lợi cho mình, và nó giám sát mọi giao dịch như một
điều kiện để tham gia.

Cái cần có là một chợ dựa trên bằng chứng mật mã thay vì niềm tin vào nền
tảng, cho phép hai bên tự nguyện bất kỳ tìm thấy nhau, thống nhất điều khoản
và xây dựng uy tín bền vững mà không cần bên thứ ba đáng tin. Trong bài này,
chúng tôi đề xuất một lời giải dùng các ý định đã ký trên relay mở, thương
lượng song phương được mã hóa, và một hệ thống uy tín trong đó bằng chứng về
giao dịch quá khứ được thiết lập bằng chữ ký số chung và được trọng số hóa
theo vị trí của mỗi người quan sát trong đồ thị giao dịch.

\section{Danh tính}

Một danh tính là một cặp khóa secp256k1, sinh ra cục bộ ngay lần dùng đầu
tiên, như trong Nostr~[2]. Không có đăng ký, không có tài khoản, không có cơ
quan cấp phát. Khóa công khai định danh người tham gia trên mọi chợ; khóa bí
mật ký mọi sự kiện người đó công bố và giải mã các tin nhắn gửi tới họ. Danh
tính có tính di động ngay từ thiết kế: uy tín tích lũy vào khóa, không vào bất
kỳ dịch vụ nào, và khóa có thể được sao lưu (mã hóa bằng cụm mật khẩu theo
NIP-49) rồi khôi phục trên bất kỳ ứng dụng nào.

Khóa là miễn phí, khiến việc tạo danh tính không tốn kém --- vấn đề trung tâm
mà thiết kế này phải xử lý. Chúng tôi sẽ quay lại nó ở các phần~6--8.

\section{Ý định}

Một bài đăng trên chợ là một \emph{ý định}: một sự kiện đã ký, có địa chỉ,
tuyên bố rằng tác giả của nó muốn mua hoặc bán thứ gì đó, ở đâu, khi nào, và
với điều khoản gì.

\begin{code}
kind:  32101 (offer) | 32102 (request)
tags:  d = stable intent id          (replaceable by author)
       t = market topic              (e.g. "sg-rideshare")
       g = geohash prefix            (location scope)
       expiration                    (NIP-40)
content: { schema, title, payload, window, flex_minutes, ... }
\end{code}

Trường \texttt{payload} được khóa theo lược đồ (\texttt{rideshare/1},
\texttt{service/1}, \dots) nên giao thức độc lập với lĩnh vực: một chợ gọi xe,
một chợ thợ sửa ống nước và một chợ đồ ăn nhà làm chỉ khác nhau ở lược đồ
payload và thẻ chủ đề. Bất kỳ ai cũng có thể định nghĩa một chợ mới bằng cách
công bố ý định dưới một thẻ mới; không tồn tại sổ đăng ký nào và cũng không
cần. Thẻ geohash giới hạn việc khám phá trong phạm vi một khu phố; thẻ hết
hạn cho phép relay dọn dẹp các bài đăng cũ. Vì ý định có địa chỉ, tác giả có
thể sửa hoặc rút một bài đăng bằng cách công bố lại dưới cùng thẻ \texttt{d}
--- việc rút chỉ là một sự kiện được công bố lại với thời hạn hết hạn gần, nên
relay loại bỏ nó đúng lịch thay vì tiếp tục phục vụ một bài đăng mà tác giả
đã từ bỏ.

Mỗi ý định mang theo một lượng nhỏ bằng chứng-công-việc (proof-of-work,
NIP-13) để việc đăng vào một chợ có một chi phí khác không, kìm hãm việc đăng
tràn lan mà không cần người gác cổng. Khối lượng tính toán là khiêm tốn ---
vài bit 0 dẫn đầu --- và ứng dụng tính nó ngoài luồng chính để giao diện vẫn
mượt trong khi tính.

Relay là hạ tầng duy nhất, và chúng là hạ tầng hàng hóa có thể thay thế cho
nhau: những kho lưu trữ đơn giản chấp nhận các sự kiện đã ký và phục vụ các
đăng ký theo bộ lọc~[2]. Một chợ tồn tại trên bất kỳ relay nào mà người tham
gia chọn dùng; sự kiểm duyệt của một relay bị đi vòng qua bằng cách công bố
lên relay khác. Ứng dụng công bố tới nhiều relay và đọc từ nhiều relay, nên
không relay nào là điểm kiểm soát.

\section{Thương lượng}

Khám phá thì công khai; thương lượng thì riêng tư. Khi một đối tác phản hồi
một ý định, hai bên trao đổi tin nhắn trực tiếp được mã hóa (NIP-04, dự kiến
NIP-17) mang các phong bì JSON điều khiển một máy trạng thái nhỏ:

\begin{code}
open --counter--> open ... (<= 8 rounds)
  |
accept --> accepted_by_us / accepted_by_them
  |                          |
  +----- second accept ------+--> confirmed
  |                                  |
  |               owner cancel ------+ (filled) --> cancelled
cancel / timeout --> cancelled / expired
\end{code}

Mỗi \texttt{counter} đề xuất một bộ điều khoản thay thế hoàn chỉnh (khung thời
gian, thanh toán, lộ trình hoặc phạm vi), nên không có sự mơ hồ về điều gì
đang được đặt lên bàn. Một giao dịch cần \emph{chấp nhận kép}: một bên chấp
nhận các điều khoản đang đứng, và giao dịch chỉ được xác nhận khi bên kia chấp
nhận lại. Điều này khép lại cuộc đua trong đó cả hai bên cùng counter đồng
thời và mỗi bên tin một phiên bản khác nhau đã được đồng ý. Giới hạn số vòng
chặn việc thương lượng spam.

Một trường hợp đặc biệt thường gặp thu gọn chấp nhận kép thành một lần chạm.
Khi một bên phản hồi chấp nhận các điều khoản đang đứng \emph{nguyên vẹn} ---
một tài xế nhận chuyến ở mức giá đã đăng --- không có điều khoản mới nào được
đưa ra, nên ứng dụng ghi nhận xác nhận ngay lập tức thay vì mở thêm một vòng.
Tuy vậy, ở phía người phản hồi, sự xác nhận vẫn là tạm thời, bởi một yêu cầu
có thể thu hút nhiều người phản hồi cạnh tranh. Chủ ý định chỉ xác nhận nhiều
nhất một người: khi làm vậy, ứng dụng của chủ rút ý định và gửi một thông báo
hủy có thẩm quyền --- \emph{đã có người nhận, chuyến đã được nhận bởi một đề
nghị khác} --- tới những người phản hồi thua cuộc. Khi nhận được, nó chuyển
một thương lượng mà người phản hồi đã cục bộ đánh dấu \texttt{confirmed} sang
\texttt{cancelled}, giải quyết cuộc đua để không hai bên nào cùng tin rằng họ
đang nắm cùng một giao dịch. Một giao dịch \texttt{completed} (đã hoàn tất) là
trạng thái cuối và không bao giờ bị đảo ngược bởi một thông báo như vậy.

Khi xác nhận, hai bên trao đổi thông tin liên hệ bên trong kênh mã hóa --- theo
thiết kế, nhiệm vụ của giao thức kết thúc ở cái bắt tay; chuyến xe hay việc
sửa chữa diễn ra ngoài đời thực.

\section{Biên nhận giao dịch}

Bản ghi thương lượng là riêng tư và một trong hai bên đều có thể bịa ra, nên
các giao dịch đã xác nhận phải được neo công khai. Khi một thương lượng đạt
trạng thái \texttt{confirmed}, mỗi bên công bố một \emph{biên nhận}:

\begin{code}
kind:  32104
tags:  d = negotiation id, p = counterparty
\end{code}

Một giao dịch được \textbf{chứng minh} khi và chỉ khi cả hai nửa cùng tồn
tại: $A$ đã ký một biên nhận nêu tên $B$, và $B$ đã ký một biên nhận nêu tên
$A$, cho cùng một negotiation id. Một khóa đơn lẻ không thể chế tạo ra một
giao dịch được chứng minh; cần hai chữ ký từ hai khóa. Điều này không ngăn
được một người kiểm soát cả hai khóa --- không sơ đồ mật mã nào ngăn được ---
nhưng nó buộc mọi giao dịch bịa đặt phải dựng nên một cạnh trong một đồ thị
công khai, và phần~\ref{sec:subjective} cho thấy vì sao các cạnh bịa đặt là vô
giá trị.

Biên nhận có địa chỉ theo negotiation id, nên mỗi bên chỉ có thể công bố nhiều
nhất một biên nhận cho mỗi giao dịch, và chúng tham chiếu đến ý định để toàn
bộ chuỗi --- bài đăng, negotiation id, xác nhận chung --- có thể được bất kỳ
ai kiểm chứng.

\section{Đánh giá}

Sau một giao dịch được chứng minh, mỗi bên có thể đánh giá bên kia một lần:

\begin{code}
kind:  32103
tags:  d = negotiation id, p = ratee
content: { score in {-1, 0, +1, +2}, note?,
           contact_verified?, contact_masked? }
\end{code}

Thẻ d làm cho đánh giá có tính lũy đẳng theo (người đánh giá, giao dịch) ---
đánh giá lại sẽ thay thế chứ không cộng dồn. Người kiểm chứng chỉ tính một
đánh giá nếu (a)~negotiation id khớp với một cặp biên nhận đã chứng minh và
(b)~người đánh giá đúng là đối tác của chính giao dịch đó. Các đánh giá tách
rời khỏi giao dịch đã chứng minh đều bị loại bỏ. Các đánh giá lặp lại giữa
cùng một cặp suy giảm theo cấp số nhân ($1, \tfrac{1}{2}, \tfrac{1}{4},
\dots$): đại lượng có ý nghĩa là \emph{số đối tác khác nhau}, đúng là đại
lượng được phần~\ref{sec:subjective} làm cho trở nên đắt đỏ.

Trường tùy chọn \texttt{contact\_verified} là một chứng thực ngang hàng rằng
người đánh giá thực sự đã liên lạc được với đối tác qua số điện thoại (đã che)
mà họ công bố, ràng buộc tuyên bố liên hệ công khai với một trao đổi riêng tư
--- xem phần~\ref{sec:privacy}. Các sự kiện đánh giá còn mang thêm một lượng
nhỏ bằng chứng-công-việc (NIP-13) làm sàn chống spam.

\section{Uy tín như một đại lượng chủ quan}
\label{sec:subjective}

Không thể xây dựng một điểm uy tín toàn cục chống được Sybil khi danh tính là
miễn phí~[4]. Thay vì chống lại định lý này, giao thức từ bỏ điểm số toàn cục.
Uy tín trong Freeport là \emph{chủ quan}: mỗi người xem tự tính trọng lượng
của từng người đánh giá từ vị trí của chính mình bằng cách duyệt theo chiều
rộng đồ thị biên nhận đã chứng minh:

\begin{center}
\begin{tabular}{lll}
\toprule
hop 0 & những người mà người xem có giao dịch đã chứng minh & trọng số 1.0 \\
hop 1 & các đối tác đã chứng minh của họ                    & trọng số 0.3 \\
hop 2 & thêm một vòng nữa                                   & trọng số 0.1 \\
---   & không thể tới được                                  & trọng số $\approx 0$ \\
\bottomrule
\end{tabular}
\end{center}

Danh sách theo dõi xã hội của người xem (NIP-02) gieo mầm cho bản đồ này để
những người mới chưa có lịch sử giao dịch vẫn nhìn qua con mắt của những người
họ đã chọn theo dõi. Uy tín hiển thị của một chủ thể khi đó là trung bình có
trọng số niềm tin của các đánh giá hợp lệ, đặt cạnh các con số thô:
\emph{số giao dịch, số đối tác khác nhau, số đối tác trong mạng lưới của bạn}.

Hãy xét một kẻ tấn công bịa ra $n$ danh tính bù nhìn và cho chúng giao dịch
rồi đánh giá lẫn nhau vô hạn. Mọi đánh giá đều qua được phép kiểm biên nhận
--- kẻ tấn công kiểm soát cả hai chữ ký. Nhưng cụm bù nhìn chỉ kết nối với một
người xem trung thực nếu có một người trung thực nào đó có một giao dịch đã
chứng minh đi vào trong cụm. Khi không có cạnh như vậy, mọi người đánh giá bù
nhìn đều mang trọng số sàn của người-không-tới-được $\varepsilon$, và ảnh
hưởng của kẻ tấn công lên nhận thức của người xem bị chặn bởi $\varepsilon$
bất kể $n$ lớn bao nhiêu. Cuộc tấn công lộ ra trong giao diện bằng chính dấu
hiệu của nó:

\begin{code}
honest:  * 12 deals - 9 partners - 3 in your network
           - verified by 9
sybil:   50 deals - 50 partners - 0 in your network
           - new account
\end{code}

Để có ảnh hưởng thực sự, kẻ tấn công phải hoàn tất những giao dịch thật với
những người tham gia thực sự có kết nối --- đến lúc đó thì hắn chỉ đơn giản là
một người tham gia. Chi phí tấn công hệ thống hội tụ về chi phí sử dụng nó một
cách trung thực, đúng là tính chất ta cần. Các bộ giảm chấn thu hẹp những
khoảng cách còn lại: một người đánh giá mà toàn bộ lịch sử biên nhận chỉ là
giao dịch với một chủ thể duy nhất bị chiết khấu (mẫu hình hai-bù-nhìn), các
tài khoản không có lịch sử sự kiện cũ hơn một tuần bị gắn cờ, và trọng số
$\varepsilon$ của người đánh giá xa lạ giữ cho các đánh giá ở xa vẫn hiển thị
nhưng không có sức ảnh hưởng.

\section{Quyền riêng tư và ràng buộc liên hệ}
\label{sec:privacy}

Danh tính mặc định là bút danh; giao thức chỉ thêm tiết lộ ở những chỗ nó đổi
được lấy niềm tin, và để người dùng tự chọn liều lượng.

\textbf{Số điện thoại.} Người tham gia có thể gắn một số điện thoại vào hồ sơ,
công bố hoặc đầy đủ (tự nguyện rõ ràng, kèm cảnh báo về tính vĩnh viễn) hoặc
che bớt chỉ còn mã quốc gia và các chữ số cuối
(\texttt{+84......678}). Việc che được áp dụng \emph{trước} khi công bố; số
đầy đủ không bao giờ rời khỏi thiết bị trừ khi ở bên trong kênh mã hóa lúc xác
nhận giao dịch. Một mặt nạ tự công bố không chứng minh được gì --- ai cũng có
thể công bố một chuỗi bất kỳ --- nên giao thức ràng buộc nó với thực tế thông
qua những người ngang hàng: lúc giao dịch, số đầy đủ truyền tới đối tác qua
kênh mã hóa; khi đánh giá, ứng dụng của đối tác chứng thực mặt nạ chuẩn tắc
của chính số mà nó \emph{đã thực sự liên lạc được}. Người kiểm chứng đối chiếu
chéo các mặt nạ được chứng thực với mặt nạ đã công bố của chủ thể và loại bỏ
các chứng thực khi không khớp. Một mặt nạ công khai giả mạo do đó bị lộ ngay
ở giao dịch trung thực đầu tiên, và ``được xác minh bởi $N$ đối tác'' nghĩa là
$N$ đối tác đã chứng minh độc lập xác nhận khả năng liên lạc --- một sự xác
minh do người ngang hàng tạo ra, không cần nhà cung cấp SMS và không cần bên
xác minh trung tâm.

\textbf{Các cuộc thương lượng} được mã hóa đầu-cuối và vô hình với relay ngoài
phần siêu dữ liệu lưu lượng. \textbf{Vị trí} chỉ được công bố ở độ chính xác
geohash do ứng dụng chọn ($\approx \pm 0.6$\,km), trung thực về chính sự thiếu
chính xác của nó. Uy tín ràng buộc vào khóa, không vào số điện thoại: việc đổi
số không chuyển đi mà cũng không thoát khỏi một lịch sử đánh giá.

\section{Động lực}

Giao thức không có token, không phí và không có người vận hành, điều này đặt
ra câu hỏi vì sao hạ tầng lại tồn tại. Câu trả lời là hạ tầng gần như miễn
phí. Relay là các kho sự kiện hàng hóa, vốn đã được hệ sinh thái Nostr vận
hành hàng trăm cái cho lưu lượng mạng xã hội; một chợ cộng đồng nằm gọn trong
bài toán kinh tế của chúng, và bất kỳ cộng đồng nào cũng có thể tự chạy relay
riêng với chi phí của một máy chủ nhỏ. Ứng dụng là những phần mềm thông thường
cạnh tranh nhau về chất lượng, không cái nào được giao thức ưu ái. Hai bên của
một giao dịch thanh toán bằng bất cứ phương tiện nào họ đã tin dùng --- tiền
mặt, chuyển khoản ngân hàng, hoặc, trong một giai đoạn giao thức về sau,
Lightning~[5], mà cơ chế ký quỹ bằng hold-invoice của nó còn có thể đóng vai
trò một tín hiệu khó giả mạo, củng cố thêm cho đồ thị ở
phần~\ref{sec:subjective}.

Điều mà thiết kế cố tình loại bỏ là tô của nền tảng: không có vị trí nào trong
giao thức để từ đó lấy một phần của các giao dịch, bởi không có nút thắt cổ
chai nào mà giao dịch buộc phải đi qua. Khám phá, thương lượng, chứng minh và
uy tín đều là các phép tính phía ứng dụng trên dữ liệu công khai, đã ký.

\section{So sánh với các nền tảng}

\begin{center}
\begin{tabular}{lll}
\toprule
 & \textbf{Nền tảng} & \textbf{Freeport} \\
\midrule
Đăng tin    & cần tài khoản, điều khoản dịch vụ & sự kiện đã ký, không cần phép \\
Ghép cặp    & thuật toán độc quyền              & bộ lọc mở, phía ứng dụng \\
Phí         & hoa hồng 15--30\%                 & không \\
Uy tín      & do nền tảng giữ, không di động    & đã ký, di động, trọng số chủ quan \\
Tranh chấp  & nền tảng phân xử                  & hệ quả về uy tín \\
Kiểm duyệt  & khóa tài khoản                    & công bố lên relay khác \\
Riêng tư    & giám sát toàn bộ giao dịch        & thương lượng mã hóa, bút danh \\
\bottomrule
\end{tabular}
\end{center}

Sự đánh đổi thành thật là khâu giải quyết tranh chấp: một nền tảng có thể hoàn
tiền cho người mua và cấm người bán; Freeport chỉ có thể bảo đảm rằng giao
dịch tồi có thể quy trách nhiệm một cách chứng minh được và hiển thị vĩnh viễn
với mọi đối tác tương lai chịu nhìn vào. Với thương mại địa phương giá trị
thấp, tần suất cao --- xe cộ, sửa chữa, đồ ăn --- hệ quả về uy tín là cơ chế
vốn đã điều chỉnh hành vi ngoài đời, và đóng góp của giao thức là làm cho cơ
chế đó trở nên di động, kiểm chứng được và không thể giả mạo ở quy mô lớn.

\section{Kết luận}

Chúng tôi đã đề xuất một chợ không có bên thứ ba đáng tin. Người tham gia
tuyên bố ý định trên relay mở, thương lượng riêng tư dưới một máy trạng thái
có giới hạn, và neo các giao dịch đã hoàn tất bằng chữ ký chung. Uy tín được
đúc lại từ một điểm số toàn cục --- bất khả thi để phòng thủ khi danh tính
miễn phí --- thành một đại lượng chủ quan tính trên đồ thị giao dịch đã chứng
minh từ chính vị trí của mỗi người xem, làm cho các cụm Sybil trở nên hiển
hiện mà không trọng lượng trong khi lịch sử trung thực thì tích lũy. Các tuyên
bố liên hệ được ràng buộc với thực tế bằng chứng thực ngang hàng thay vì bằng
bất kỳ bên xác minh nào. Giao thức nhỏ gọn: bốn loại sự kiện và một phong bì
tin nhắn trên một mạng relay sẵn có. Bản thân các chợ là những quy ước --- một
thẻ và một lược đồ --- nên ai cũng có thể mở một chợ, và không ai có thể đóng
nó lại.

\begin{thebibliography}{9}
\small

\bibitem{nakamoto} S.~Nakamoto, ``Bitcoin: A Peer-to-Peer Electronic Cash
System,'' 2008.

\bibitem{nostr} fiatjaf et al., ``Nostr: Notes and Other Stuff Transmitted
by Relays --- NIP-01: Basic protocol flow,''
\texttt{github.com/nostr-protocol/nips}.

\bibitem{nips} Các NIP của Nostr được dùng: NIP-02 (danh sách theo dõi),
NIP-04/NIP-17 (nhắn tin mã hóa), NIP-13 (bằng chứng công việc), NIP-19 (thực
thể bech32), NIP-40 (hết hạn), NIP-44 (mã hóa có phiên bản), NIP-49 (mã hóa
khóa), NIP-96/98 (lưu trữ tệp HTTP và xác thực).

\bibitem{douceur} J.~R.~Douceur, ``The Sybil Attack,'' \emph{IPTPS}, 2002.

\bibitem{lightning} J.~Poon, T.~Dryja, ``The Bitcoin Lightning Network:
Scalable Off-Chain Instant Payments,'' 2016.

\bibitem{sybilrank} Q.~Cao, M.~Sirivianos, X.~Yang, T.~Pregueiro, ``Aiding
the Detection of Fake Accounts in Large Scale Social Online Services''
(SybilRank), \emph{NSDI}, 2012.

\bibitem{sybilguard} H.~Yu, M.~Kaminsky, P.~B.~Gibbons, A.~Flaxman,
``SybilGuard: Defending Against Sybil Attacks via Social Networks,''
\emph{SIGCOMM}, 2006.

\end{thebibliography}

\end{document}
